% ============================================================
% Chapter 6: Experimental Design
%   Section: Construction of the Custom Validation Set
% Included from chapters/06-experimental-design/experimental-design.tex
% ============================================================

% Follows ../02-evaluation-dataset/, whose subsection "Domain
% Limitation and a Custom Validation Set" (\ref{subsec:custom-gold-set})
% motivates this set. Requires booktabs, tabularx and natbib in the
% main preamble; its BibTeX entries are in references.bib.

\section{Construction of the Custom Validation Set}
\label{sec:custom-dataset}

Section~\ref{subsec:custom-gold-set} argued that X-Fact
\cite{gupta2021xfact}, although methodologically aligned with the
open-retrieval architecture of the proposed method, does not cover
its application domain: X-Fact is dominated by political statements,
whereas the system is applied to positive-news articles about
science, the environment, health, culture and society, whose claims
mix checkable facts with interpretive readings of them. This section
describes how the complementary, hand-labelled validation set was
built: its size and stratification, the record format, the label
vocabulary, the annotation procedure, the tie-break rules fixed
\emph{before} labelling began, and the protocol used to verify the
consistency of the labels.

% ============================================================
% Chapter 6: Experimental Design
%   Section: Construction of the Custom Validation Set
%     Subsection: Size and Stratification
% Included from chapters/06-experimental-design/03-custom-validation-set/custom-validation-set.tex
% ============================================================

\subsection{Size and Stratification}
\label{subsec:custom-size}

The set targets \textbf{150 claims}, with 100 as the minimum below
which per-stratum counts stop being informative. With 150 claims, a
95\% confidence interval on an accuracy near 50\% is approximately
$\pm 8$ percentage points ($1.96\sqrt{0.25/150}\approx0.080$), which
is sufficient to separate the models compared in
Chapter~\ref{ch:results} when they differ substantially, and
insufficient to rank models that differ by a few points. That limit
is stated here rather than discovered later.

The claims are stratified along two axes, so that neither a single
verdict nor a single subject dominates the metrics:

\begin{itemize}
    \item \textbf{Verdict}: the five labels of the system's own
    \texttt{Verdict} enumeration (Table~\ref{tab:custom-labels}).
    \item \textbf{Topic}: each claim carries one of the 23 topics of
    the system's topic classifier (\texttt{src/config/topics.py}).
    Because 23 topics over 150 claims leaves six or seven claims per
    topic, too few to balance on, balance is sought over the five
    \emph{topic groups} into which those topics are already
    organised: society, science, environment, culture and health.
\end{itemize}

Five verdicts by five groups gives 25 cells, with a target of six
claims each (30 per verdict and 30 per group). Language (English or
Spanish) and claim type (factual, numerical or interpretive) are
recorded and reported, but not targeted.

The targets are a goal, not a quota. Positive-news articles are
predominantly accurate, so the \texttt{FALSE} and \texttt{MISLEADING}
cells are expected to be the hardest to fill from the outlet's own
articles. A cell that remains short is reported with its real count
rather than padded with claims from outside the domain, and the
evaluation reports macro-averaged $F_1$ alongside accuracy for that
reason: under class imbalance, accuracy rewards a system that
predicts the majority verdict.


% ============================================================
% Chapter 6: Experimental Design
%   Section: Construction of the Custom Validation Set
%     Subsection: Record Format
% Included from chapters/06-experimental-design/03-custom-validation-set/custom-validation-set.tex
% ============================================================

\subsection{Record Format}
\label{subsec:custom-format}

Every record reproduces, key for key and in the same order, the
format into which X-Fact was converted for this work
(Section~\ref{subsec:xfact}), so that the harness that scores the
system against X-Fact scores it against this set without
modification. The shared keys are listed in
Table~\ref{tab:custom-format}; the custom set appends its own keys
after them, which a reader of the X-Fact format ignores.

\begin{table}[ht]
\centering
\small
\begin{tabularx}{\linewidth}{@{}lX@{}}
\toprule
\textbf{Key} & \textbf{Content in the custom set} \\
\midrule
\multicolumn{2}{@{}l}{\emph{Shared with X-Fact}} \\
\texttt{language} & \texttt{en} or \texttt{es}. \\
\texttt{site} & Domain of the outlet in which the claim was read
(X-Fact records the fact-checking organisation here). \\
\texttt{claimant} & Who makes the claim: a person, organisation or
the outlet itself. \\
\texttt{claim} & One atomic, checkable assertion, as the article
states it. \\
\texttt{claimDate} & Publication date of the article. \\
\texttt{reviewDate} & When the label was last set. \\
\texttt{labelRaw} & The label in the annotation guide's vocabulary
(Table~\ref{tab:custom-labels}). \\
\texttt{label} & The label in the system's \texttt{Verdict}
vocabulary. \\
\texttt{referenceEvidenceLinks} & URLs of the evidence on which the
verdict rests. \\
\texttt{split} & Always \texttt{test}: no claim is ever used to
train or tune on. \\
\midrule
\multicolumn{2}{@{}l}{\emph{Added by the custom set}} \\
\texttt{id} & The record's file name (\texttt{fact001},
\texttt{fact002}, \ldots), in the order of annotation. \\
\texttt{topic} & One of the 23 classifier topics; its group is
derived. \\
\texttt{claimType} & \texttt{factual}, \texttt{numerical} or
\texttt{interpretive}. \\
\texttt{sourceTier} & Tier of the strongest evidence used
(Rule~4, Section~\ref{subsec:custom-rules}). \\
\texttt{onlyOwnSource} & Whether only the organisation the story is
about backs the claim (Rule~1). \\
\texttt{evidenceDate} & Date of the newest evidence used
(Rule~3). \\
\texttt{articleUrl} & The article the claim was taken from. \\
\texttt{annotatorNote} & The reasoning behind a close call, and the
rule that decided it. \\
\texttt{review} & The blind second label, when the record is in the
verification sample (Section~\ref{subsec:custom-review}). \\
\bottomrule
\end{tabularx}
\caption{Record format of the custom validation set. The first ten
keys are those of the X-Fact records used in this work, in the same
order.}
\label{tab:custom-format}
\end{table}

As in X-Fact, the evidence links are kept for inspecting a
disagreement, not fed to the system: the system retrieves its own
evidence from the open web at verification time, and handing it the
annotator's sources would measure the reasoning stage alone, which is
precisely the flaw for which FEVER was discarded
(Section~\ref{subsec:discard-justification}).


% ============================================================
% Chapter 6: Experimental Design
%   Section: Construction of the Custom Validation Set
%     Subsection: Label Vocabulary
% Included from chapters/06-experimental-design/03-custom-validation-set/custom-validation-set.tex
% ============================================================

\subsection{Label Vocabulary}
\label{subsec:custom-labels}

X-Fact preserves each fact-checker's original label
(\texttt{labelRaw}) next to a normalised one (\texttt{label}); the
custom set does the same. The annotation guide names its labels after
AVeriTeC \cite{schlichtkrull2023averitec}, the real-world,
open-web claim verification benchmark whose four classes ---
\emph{Supported}, \emph{Refuted}, \emph{Not Enough Evidence} and
\emph{Conflicting Evidence/Cherry-picking} --- are defined by the
evidence available rather than by an absolute notion of truth, which
is the stance this work also takes. A fifth label,
\emph{partially supported}, is added because the system's own
\texttt{Verdict} distinguishes a claim whose central assertion holds
while a detail does not, the outcome in which real verification most
often lands. Table~\ref{tab:custom-labels} gives the correspondence.

\begin{table}[ht]
\centering
\small
\begin{tabularx}{\linewidth}{@{}llX@{}}
\toprule
\textbf{\texttt{labelRaw}} & \textbf{\texttt{label}} &
\textbf{Definition} \\
\midrule
supported & \texttt{TRUE} & The evidence supports the claim as
stated. \\
partially supported & \texttt{PARTIALLY\_TRUE} & The central claim
holds; a detail does not. \\
conflicting evidence/cherrypicking & \texttt{MISLEADING} & Sources
of equal standing conflict, or the claim is true but selected to
mislead. \\
refuted & \texttt{FALSE} & The evidence refutes the claim. \\
not enough evidence & \texttt{UNVERIFIED} & There is not enough
independent evidence to decide either way. \\
\bottomrule
\end{tabularx}
\caption{Annotation labels, after AVeriTeC
\cite{schlichtkrull2023averitec}, and the system verdicts they map
to. The X-Fact mapping used in this work sends X-Fact's
\emph{partly true/misleading} to \texttt{MISLEADING} and its
\emph{complicated/hard to categorise} to \texttt{UNVERIFIED}, so both
sets are scored on the same five classes.}
\label{tab:custom-labels}
\end{table}


% ============================================================
% Chapter 6: Experimental Design
%   Section: Construction of the Custom Validation Set
%     Subsection: Annotation Procedure
% Included from chapters/06-experimental-design/03-custom-validation-set/custom-validation-set.tex
% ============================================================

\subsection{Annotation Procedure}
\label{subsec:custom-procedure}

Claims are taken from articles published by the positive-news and
science sources the system ingests, and neither the articles nor the
claims are chosen by the annotator. Each day a batch of about ten
articles is drawn at random from the sources' current feeds and topic
pages: at most one article per source, the two languages alternated,
every article already labelled or previously proposed excluded, and
articles whose predicted topic belongs to an under-filled topic group
preferred, after articles published within the previous sixty days,
which keep small the gap between the evidence admissible under
Rule~3 and the present-day web the system searches. The draw is seeded
with the date, so the same set of candidate links yields the same
draw.
For each article, the system's own claim extractor and anchor-claim
selector propose between one and three claims, which are therefore
the claims the system would itself verify.

The annotator either labels a proposed claim or rejects it with a
recorded reason: an opinion or interpretation, a prediction, a
checkable but trivial statement, a fragment of a sentence, or a
duplicate. Those rejections are kept, and the share of proposals
accepted is reported as the precision of the system's claim
selection as judged by the annotator. Each record holds one atomic
claim, kept in the article's own wording; a proposed sentence that
asserts two things is split into two records, following the claim
decomposition practice of real-world fact-checking datasets
\cite{guo2022survey}. The annotator then searches for evidence
independently of the system, records the links on which the verdict
rests, and assigns the label under the rules of
Section~\ref{subsec:custom-rules}. The system's verdict is never
shown during annotation.

Records are entered through a small web form built for this purpose
and kept apart from the system itself: a single program using only the
Python standard library, so that labelling depends neither on the
system's services nor on any installed package. Its fields are exactly
the keys of Table~\ref{tab:custom-format}, and it enforces the rules
that can be checked mechanically at the moment a record is saved:
a verdict other than \texttt{UNVERIFIED} requires at least one
evidence link; a claim backed only by its protagonist's own source
must be labelled \texttt{UNVERIFIED} and the source named; and an
evidence date later than the article's publication date is rejected.
The same form shows the running count per verdict and topic group
against the targets of Section~\ref{subsec:custom-size}, which is how
balance is steered during annotation rather than assessed after it.

Each record is stored as its own file (\texttt{fact001.json},
\texttt{fact002.json}, \ldots) in the project's version-controlled
repository. Every verified claim is therefore individually visible,
and every later correction is recorded as a change to that one file,
which leaves an audit trail of the annotation. Once the set is
complete, the files are validated again and merged, in order, into a
single JSON Lines file in the X-Fact format, which is what the
evaluation harness reads. The per-record files remain the source, and
the merged file is derived from them.


% ============================================================
% Chapter 6: Experimental Design
%   Section: Construction of the Custom Validation Set
%     Subsection: Tie-Break Rules
% Included from chapters/06-experimental-design/03-custom-validation-set/custom-validation-set.tex
% ============================================================

\subsection{Tie-Break Rules}
\label{subsec:custom-rules}

The decisions that cause the most hesitation during annotation are
also those on which a single annotator is least consistent with
themself over time. The following four rules were therefore fixed
before labelling began and are applied without exception; the
\texttt{annotatorNote} records which rule decided a close call.

\paragraph{Rule 1 --- The organisation's own source.}
Positive-news claims frequently originate in the press release of
the very NGO, company or institution the article is about. A claim
for which the \emph{only} available evidence is that organisation's
own communication (press release, report, website or social media
post) is labelled \emph{not enough evidence}
(\texttt{UNVERIFIED}), and the source is named in the note. An
organisation reporting on itself is not independent confirmation,
and a claim the system could only confirm by reading the claimant's
own release is not one it has verified.

\paragraph{Rule 2 --- Numeric tolerance.}
A figure is considered supported when it lies within $\pm10\%$ in
relative terms of the figure in the source (an article's ``30\%''
against a source's ``28\%'' is supported), or when it is a reasonable
rounding of it (``nearly 3,000'' for 2,870). Beyond that margin, the
claim is \emph{partially supported} if the figure is a detail of an
otherwise correct claim, and \emph{refuted} if the figure is the
claim.

\paragraph{Rule 3 --- Date of the evidence.}
Only evidence available on the article's publication date is
admissible, as in AVeriTeC \cite{schlichtkrull2023averitec}, which
restricts evidence to documents published before the claim to avoid
temporal leakage. A claim that could not be checked when it was
published and was confirmed later remains \emph{not enough evidence};
the verdict does not change. Evidence that did not yet exist cannot
have made a claim well-founded when it was made, and fact-checking
evaluations that ignore this condition have been shown to overstate
what a system could do in practice \cite{glockner2022missing}.

\paragraph{Rule 4 --- Hierarchy of sources.}
Sources are ranked as follows: primary source (a report, official
statistic or peer-reviewed study) $>$ reference media outlet $>$
press release $>$ social media. When sources conflict, the higher
tier prevails; when sources of the same tier conflict and none is
decisive, the claim is labelled \emph{conflicting evidence}
(\texttt{MISLEADING}). The tier of the strongest source used is
recorded in \texttt{sourceTier}, which also allows the evaluation to
be broken down by how well-sourced each claim was.


% ============================================================
% Chapter 6: Experimental Design
%   Section: Construction of the Custom Validation Set
%     Subsection: Verification of the Labels
% Included from chapters/06-experimental-design/03-custom-validation-set/custom-validation-set.tex
% ============================================================

\subsection{Verification of the Labels}
\label{subsec:custom-review}

The set is labelled by a single annotator, so inter-annotator
agreement cannot be measured. What can be measured is
\textbf{intra-annotator agreement}: whether the guide, applied by the
same person at a different time, produces the same label. Once
labelling is complete, 20\% of the records are labelled a second
time:

\begin{itemize}
    \item \textbf{The sample is not chosen by the annotator.} It is
    the first $\lceil 0.2\,n \rceil$ records in the order of a
    cryptographic hash (SHA-256) of their identifier, which is fixed
    in advance and independent of content.
    \item \textbf{The second pass is blind.} The review view shows
    the claim, its source article, dates, topic and evidence links,
    but not the first label, the \texttt{onlyOwnSource} flag or the
    annotator's note, all of which reveal it.
    \item \textbf{It takes place at least a week after the first
    pass}, so that the second label is produced by the guide rather
    than recalled.
\end{itemize}

Agreement is reported as observed agreement and as Cohen's $\kappa$
\cite{cohen1960coefficient}, which discounts the agreement two
labellings would reach by chance given how often each used every
label; this matters here because the label distribution is expected
to be uneven. Following the conventions reviewed by Artstein and
Poesio \cite{artstein2008intercoder}, $\kappa \geq 0.6$ (``substantial''
on the scale of Landis and Koch \cite{landis1977measurement}) is taken
as the threshold for the guide to be considered consistently
applicable; below it, the rules behind the disagreements are revised
and the affected strata relabelled.

Disagreements are resolved against the guide and, where the first
label proves wrong, the record is corrected. Agreement remains
measured on the \emph{first} label, which the record keeps beside the
second: correcting a record after its review must not turn a
disagreement into an agreement.


% ============================================================
% Chapter 6: Experimental Design
%   Section: Construction of the Custom Validation Set
%     Subsection: Limitations
% Included from chapters/06-experimental-design/03-custom-validation-set/custom-validation-set.tex
% ============================================================

\subsection{Limitations}
\label{subsec:custom-limitations}

\begin{itemize}
    \item \textbf{One annotator.} Self-agreement measures the
    consistency of the guide's application, not its validity: a
    systematic misreading applied twice agrees with itself.
    \item \textbf{Selection by the system.} Articles are drawn at
    random and claims are proposed by the system, which removes the
    annotator's tendency to pick the most checkable cases, but ties
    the set to the claims the system's selector favours. A claim the
    selector would never propose cannot enter the set, so the
    evaluation measures verification on the system's own claim
    choice, not on an independent sample of all check-worthy claims.
    The first six records were labelled before the daily draw existed
    and were chosen by hand; they are identified as such (no batch
    refers to them) and reported separately.
    \item \textbf{Temporal asymmetry.} Rule~3 restricts the
    annotator to evidence available at publication time, whereas the
    system searches today's web, without a date restriction. A system
    verdict informed by later evidence can therefore be scored as an
    error; such cases are identified from the evidence dates and
    reported separately rather than silently absorbed into the error
    rate.
    \item \textbf{Size.} As computed in
    Section~\ref{subsec:custom-size}, differences of a few points
    between models are within the uncertainty of a 150-claim set.
\end{itemize}

